% chap2.tex
\part{구조 및 인터페이스 명세}
\label{part:ioref}

% ---------------------------------------------------------------------------- %
%                                  NEW SECTION                                 %
% ---------------------------------------------------------------------------- %

\chapter{디렉토리 구조}
\simulator 는 핵심 알고리즘을 포함하는 python 모듈(\pymodule)과 함께, 독립 실행형 (portable) Python, 실행 스크립트, 기술문서, 및 예시파일을 포함한다.
디렉토리의 각 구성 요소별 내용은 \textcref{table:release-dirstructure}\과 같다.

\begin{defaulttable}
  \caption{\simulator~디렉토리 구성}
  \label{table:release-dirstructure}
  \begin{tabularx}{\textwidth}{l X c}
    \toprule
    \textbf{구성 요소 (경로)} & \textbf{내용} & \textbf{참조} \\
    \midrule
    
    \textbf{src/} & 엔진의 핵심 알고리즘이 포함된 소스 코드 & - \\
    
    \hspace{0.5em} $\llcorner$ idragon/ & EnergyPlus 친화적 구조 및 idf 변환 & - \\
    \hspace{1.5em} $\llcorner$ \_energyPlus/ & 내장된 EnergyPlus 엔진 & \ref{section:dirstructure-devenv} \\
    \hspace{0.5em} $\llcorner$ \pymodule/ & \simulator 친화적 구조 및 중간 변환 & \ref{section:howtorun-pythonmodule} \\
    \hspace{0.5em} $\llcorner$ launcher/ & Excel 기반 입력파일의 실행을 위한 web 기반 GUI & \ref{section:howtorun-excellauncher} \\
    \addlinespace 
    
    \textbf{docs/} & 엔진 관련 명세 & \ref{section:dirstructure-documentation} \\
    \hspace{0.5em} $\llcorner$ \texttt{TechnicalReferenceManual} & \simulator 의 구조 및 알고리즘 명세 (본 문서) & - \\
    \hspace{0.5em} $\llcorner$ \texttt{RegressionTestReport} & 버전 검증용 테스트 파일 실행 결과 & - \\
    \hspace{0.5em} $\llcorner$ \texttt{ReleaseNote} & 버전 관리 규칙 및 업데이트 내용 & - \\
    \addlinespace

    \textbf{examples/} & 입출력 파일 예시 및 포맷 & - \\
    \hspace{0.5em} $\llcorner$ grm/ & 입력파일 예시 & - \\
    \hspace{0.5em} $\llcorner$ grr/ & 출력파일 예시 & - \\
    \addlinespace

    \textbf{venv/} & 필요한 라이브러리가 포함된 독립 실행형 Python & \ref{section:dirstructure-devenv} \\
    \addlinespace
    
    \texttt{runEngine} & 엔진 실행 스크립트 & \ref{section:howtorun-cli} \\
    \texttt{runExcelLauncher} & Excel 기반 입력파일의 실행을 위한 web 기반 GUI 실행 스크립트 & \ref{section:howtorun-excellauncher} \\
    
    \bottomrule
  \end{tabularx}
\end{defaulttable}

\section{외부 라이브러리}
\label{section:dirstructure-devenv}

\begin{itemize}
  \item \textbf{Python} \simulator \는 2023년 말 출시되어 현재 가장 널리 쓰이는 버전 중 하나인 \texttt{Python 3.12}\cite{python312}를 사용한다.
  \item \textbf{EnergyPlus} \simulator 에서 EnergyPlus 실행파일 변환 및 실행을 담당하는 idragon 모듈은 2024년 9월 말에 출시된 \texttt{EnergyPlus 24.2}\cite{energyplus242} 버전에 맞추어 개발되었다.
\end{itemize}


\section{기술문서}
\label{section:dirstructure-documentation}

\begin{itemize}
  \item \textbf{Technical Reference Manual} 본 문서.
  \item \textbf{Regression Test Report} Regression Test Report는 소프트웨어의 기능 추가나 코드 수정, 혹은 종속 라이브러리의 업데이트가 기존에 정상적으로 작동하던 기능에 의도치 않은 오류(Regression)를 유발하지 않았음을 검증하고 기록한 문서이다\cite{sommerville2015software}. 이 보고서에서는 표준 테스트용 시뮬레이션을 수행한 결과값, 실행 속도 등을 정량적으로 제시하므로, 이전 버전과 현재 버전 간의 차이를 비교, 엔진의 수치적 정합성과 시스템 안정성을 확인 가능하다.
  \item \textbf{Release Note} \simulator 의 버전별 호환성 등 버전 관리 규칙을 명세하고, 각 버전별 업데이트 내용을 추가, 변경, 수정, 제거 4가지 항목으로 구분하여 기술한다.
\end{itemize}


\chapter{실행 방법}

\simulator\은 아래 세 가지 방법으로 실행할 수 있다.
\begin{itemize}
  \item \textbf{명령줄 인터페이스 (Command Line Interface; CLI)}: runEngine.bat 스크립트 파일을 이용, cmd 콘솔 명령어로 \simulator 의 주요 기능을 실행한다. GUI 프로그램 등, 본 엔진과 연동되는 외부 프로그램에 권장되는 방법이다. 포맷 변환, DB 조회 등 외부 프로그램과의 연동이 필요한 기능들이 지원된다.
  \item \textbf{Excel 기반 입력파일 실행}: Excel 기반 입력파일(\texttt{.xlsx} 확장자)을 작성하여, 웹 기반 GUI를 통해 \simulator \를 실행한다. 본 엔진을 독립적으로 사용하는 사용자에게 권장되는 방법이며, 시뮬레이션의 실행만 가능하다.
  \item \textbf{Python 모듈 직접 사용}: 엔진에 포함된 portable Python, 또는 별도의 환경에서 \pymodule\을 직접 import하여 \simulator 의 모든 함수를 활용할 수 있다. 본 엔진을 parametric 시뮬레이션 등 변칙적으로 활용하려는 연구자 및 개발자에게 권장되는 방법이다.
\end{itemize}


\section{명령줄 인터페이스 (CLI)}
\label{section:howtorun-cli}

\simulator\은 명령행 인터페이스(Command-Line Interface, CLI)를 통해 실행할 수 있으며, 배포 패키지에는 이를 위한 배치 스크립트\inlinecode{runEngine.bat}가 포함되어 있다. 스크립트는 내부적으로 \simulator 에 포함된 Python을 이용, epsimple 모듈을 실행한다 즉, 아래 CLI signature는 \inlinecode{venv/python.exe -m epsimple <command> [args]}와 동일하며, 외부 연동 시 안정성을 위해 Python의 직접 호출보다는 배치 스크립트를 이용하는 것이 권장된다.

\begin{codeblock}{CLI command signature}
runEngine <command> [args]
\end{codeblock}

여기서 \inlinecode{<command>}는 시뮬레이션 실행 (\texttt{run}: \ref{subsection:clisubcommands-run}), 포맷 변환(\texttt{convert}: \ref{subsection:clisubcommands-convert}) 두 가지가 지원된다.

\subsection{시뮬레이션 실행}
\label{subsection:clisubcommands-run}

\simulator 의 핵심 기능으로, 입력파일을 받아 시뮬레이션을 실행하고, 결과를 출력파일로 저장하는 기능이다. 입력파일은 \texttt{grm}(json 형식, \ref{section:format-explanation-grm}) 또는 \texttt{xlsx}(\simulator 전용 엑셀 입력 형식, \ref{section:format-explanation-xlsx}) 모두 가능하며, 출력파일은 \texttt{grr}(json 형식, \ref{section:format-explanation-grr})이다.

\begin{codeblock}{CLI subcommand signature: `run'}
runEngine run --input <input_filepath> [--output <output_filepath>]
\end{codeblock}

\paragraph{입력 인자}
\begin{description}
  \item[\texttt{-i}, \texttt{--input}] 입력파일의 경로. \texttt{grm} 또는 \texttt{xlsx} 형식의 파일이 가능하다. 입력파일의 형식은 확장자에 따라 \texttt{xlsx} 확장자는 Excel 기반 입력파일로, 그 외는 json 형식의 입력파일로 인식된다.
  \item[\texttt{-o}, \texttt{--output}] 출력파일의 경로. 지정하지 않으면 입력파일의 경로에서 확장자만 \texttt{grr}로 바꾼 경로에 저장된다.
\end{description}

\subsection{포맷 변환}
\label{subsection:clisubcommands-convert}

입력파일의 형식을 변환하는 기능이다. \simulator 에서 입력파일은 \texttt{xlsx} $\rightarrow$ \texttt{grm} $\rightarrow$ \texttt{idf} 순으로 변환되어 시뮬레이션되므로, \texttt{xlsx} $\rightarrow$ \texttt{idf}를 포함하여 3가지 변환이 가능하다.

\begin{codeblock}{CLI subcommand signature: `convert'}
runEngine convert <input_filepath> --source <input_type> --destination <output_type> [--output_filepath <output_filepath>]
\end{codeblock}

\paragraph{입력 인자}
\begin{description}
  \item[\texttt{input\_filepath}] 변환할 입력파일의 경로. \texttt{grm} 또는 \texttt{xlsx} 형식의 파일이 가능하다.
  \item[\texttt{-s}, \texttt{--source}] 입력파일의 형식. 실제 형식에 맞게 \texttt{excel}, \texttt{json} 중 하나를 지정할 수 있다.
  \item[\texttt{-d}, \texttt{--destination}] 출력파일의 형식. \texttt{json}(입력파일의 형식이 \texttt{excel}인 경우) 또는 \texttt{idf} 중 하나를 지정할 수 있다.
  \item[\texttt{-o}, \texttt{--output\_filepath}] 출력파일의 경로. 지정하지 않으면 입력파일의 경로에서 확장자만 변경된 경로에 저장된다.
\end{description}

\section{Excel 기반 입력파일의 실행 방법}
\label{section:howtorun-excellauncher}

배치 스크립트 \inlinecode{runExcelLauncher.bat}를 실행하면, 서버 역할을 하는 command window와 함께 웹 기반 GUI가 열리며, Excel 기반 입력파일(\texttt{.xlsx} 확장자)을 업로드하여 \simulator \를 실행할 수 있다. GUI는 입력파일의 형식이 올바른지 검증한 뒤, 시뮬레이션을 실행하고, 결과를 레포트 형식으로 제공한다.

웹 기반 GUI는 최대 10개의 리모델링 이후 모델을 동시에 업로드, 시뮬레이션 및 비교할 수 있으며, 최대 4개 코어를 활용하는 병렬 시뮬레이션을 지원한다.


\section{Python 모듈(\pymodule) 직접 사용법}
\label{section:howtorun-pythonmodule}

\pymodule 의 주요 함수는 \texttt{/src/epsimple/api.py}에 정의되어있으며, \pymodule 을 import하여 사용이 가능하다. 주요 함수 및 입력 인자는 \ref{section:howtorun-cli}에서 기술한 CLI와 유사하다.

\begin{codeblock}{Python module API example}
from epsimple import (
  run_grjson         ,
  run_grexcel        ,
  convert_inputformat,
)

run_grjson(
  "my_input.grm" ,
  "my_output.grr",
)
\end{codeblock}


\chapter{입력 및 출력 파일 명세}

\section{입력파일 (\excel)}
\label{section:format-explanation-excel}

\excel 의 입력파일은 실제 건물의 물리적 구성 관계를 반영하여 계층적으로 구성되어 있다 \parencref{fig:grm_structure}. 즉, 건물(\texttt{building})을 최상위 객체로 하여 층(\texttt{floor}), 실(\texttt{zone}), 면(\texttt{surface}), 개구부(\texttt{fenestration}) 등과 같은 공간 요소가 실제 건물의 포함 관계와 유사한 형태로 단계적으로 조직된다. 이러한 구조를 통해 건물의 공간적 구성과 시스템 연결 관계를 직관적으로 표현할 수 있다.

한편, 재료(\texttt{material})나 구조체(\texttt{construction})과 같은 물리적 속성은 공간 객체 내부에 직접 중복 정의되지 않고 별도의 객체로 정의된 뒤 참조 방식으로 연결된다. 예를 들어 면(\texttt{surface})은 특정 구조체(\texttt{surface\_construction})를 참조하는 방식으로 구성된다. 이러한 구조는 동일한 정보를 여러 객체에서 재사용할 수 있도록 한다. 또한, 설비 시스템(\texttt{hvac\_system})과 제어 로직(\texttt{control\_logic})은 건물 객체의 하위에 별도의 분류로 정의되어, 공간 요소와는 독립적으로 구성되며, 필요한 경우 공간 요소와 연결될 수 있도록 설계되어 있다. 이는 건물의 물리적 구성과 시스템 구성을 명확히 구분한다.

\begin{figure}[ht]
  \includegraphics[width=0.99\textwidth]{grm-datastructure.png}
  \caption{\simulator\ 입력파일 체계도}
  \label{fig:grm_structure}
\end{figure}




\subsection{규칙}

\subsubsection{변수의 명칭 관련}

\paragraph{사용가능한 문자} 모든 변수의 명칭은 영문 소문자 및 밑줄(\_)로만 구성한다.

\subsubsection{이름 기반 객체 참조 관련}
\label{subsubsection:ioref:id_description}
\excel 에서 ``이름''의 값을 입력하는 경우, 객체가 생성되며, 다른 객체의 지칭은 ``이름''값에 기반하여 이루어진다.

\paragraph{사용가능한 문자} 모든 ``이름''은 한글 혹은 영문 대소문자, 숫자, 하이픈(-) 및 밑줄(\_)로만 구성하며, 특수 문자로는 작성하지 않는다.

\paragraph{이름의 중복} 동일한 ``이름''은 같은 \texttt{class} (\excel 에서, 각 \texttt{sheet}를 의미한다.)의 객체끼리 공유할 수 없다. 
한글 혹은 영문 대소문자, 숫자, 하이픈(-) 및 밑줄(\_)로만 구성하며, 특수 문자로는 작성하지 않는다.



\subsubsection{변수의 유효성}

% ------------------------------- 여기서만 쓰는 함수들 ------------------------------- %

% --- 공통 chip (그대로 사용) ---
\newtcbox{\tagchip}[1][]{%
  on line, boxsep=0.6ex, left=0.9ex, right=0.9ex, top=0.2ex, bottom=0.2ex,
  arc=2.2pt, boxrule=0pt, colframe=white,
  colback=black, coltext=white, fontupper=\sffamily\bfseries\footnotesize,
  #1}

% --- 타입 색상 ---
\definecolor{typeString}{RGB}{206,145,120}
\definecolor{typeFloat}{RGB}{181,206,168}
\definecolor{typeInt}{RGB}{78,201,176}
\definecolor{typeEnum}{RGB}{86,156,214}
\definecolor{typeBool}{RGB}{142,206,221}

% --- 요구수준 색상 ---
\definecolor{reqBlueDark}{RGB}{37,99,235}
\definecolor{reqBlueLight}{RGB}{147,197,253}
\definecolor{reqGrayDark}{RGB}{107,114,128}
\definecolor{reqGrayLight}{RGB}{209,213,219}


% --- 엑셀 요구수준 색상 ---
\definecolor{excelRequired}{RGB}{169,209,142}
\definecolor{excelOptional}{RGB}{155,194,230}
\definecolor{excelDisabled}{RGB}{191,191,191}


% --- TypeTag (S/F/I/E) ---
\newcommand{\TypeTag}[1]{%
  \begingroup
  \ifthenelse{\equal{#1}{S}}{\tagchip[colback=typeString,coltext=black]{S}}{%
  \ifthenelse{\equal{#1}{F}}{\tagchip[colback=typeFloat ,coltext=black]{F}}{%
  \ifthenelse{\equal{#1}{I}}{\tagchip[colback=typeInt   ,coltext=white]{I}}{%
  \ifthenelse{\equal{#1}{E}}{\tagchip[colback=typeEnum  ,coltext=white]{E}}{%
  \ifthenelse{\equal{#1}{B}}{\tagchip[colback=typeBool  ,coltext=white]{B}}{%
  \tagchip[colback=gray!60,coltext=black]{#1}}}}}}%
  \endgroup
}

% --- ReqTag (R/CR/O/CO) ---
\newcommand{\ReqTag}[1]{%
  \begingroup
  \ifthenelse{\equal{#1}{R}} {\def\shade{reqBlueDark}\def\fg{white}}{}%
  \ifthenelse{\equal{#1}{CR}}{\def\shade{reqBlueLight}\def\fg{black}}{}%
  \ifthenelse{\equal{#1}{O}} {\def\shade{reqGrayDark}\def\fg{white}}{}%
  \ifthenelse{\equal{#1}{CO}}{\def\shade{reqGrayLight}\def\fg{black}}{}%
  \ifthenelse{\isundefined{\shade}}{\def\shade{gray!60}\def\fg{black}}{}%
  \tagchip[colback=\shade,coltext=\fg]{#1}%
  \endgroup
}


% --- excel tag ---
\newcommand{\InputTag}[1]{%
  \begingroup
  \ifthenelse{\equal{#1}{R}}{\tagchip[colback=excelRequired,coltext=black]{필수입력}}{%
  \ifthenelse{\equal{#1}{O}}{\tagchip[colback=excelOptional,coltext=black]{optional}}{%
  \ifthenelse{\equal{#1}{D}}{\tagchip[colback=excelDisabled,coltext=black]{입력불가}}{%
  \tagchip[colback=gray!60,coltext=black]{#1}}}}%
  \endgroup
}


% --- table definition --- 
\newcommand{\jsontableunit}[1]{%
  \ifthenelse{\equal{#1}{}}{-}{$#1$}%
}

\newcommand{\jsontabledash}[1]{%
  \ifthenelse{\equal{#1}{}}{-}{#1}%
}

\newcommand{\jsontablerow}[7]{%
  \nolinkurl{#1} & #2 & #3 &
  \jsontabledash{#4} & \jsontabledash{#5} & \jsontabledash{#6} & \jsontableunit{#7} \\
}

\newcommand{\needtobeexplain}{상세설명 참고}

\newcommand{\jsontable}[2]{
  \begin{table}[htbp]
    \centering
    \caption{grjson 구조에서 \texttt{#1} 클래스의 속성들}
    {\footnotesize
    \begin{tabularx}{\textwidth}{l c c c Y Y c}
      \toprule
      변수명 & 타입 & 필수 & 기본값 & 조건 & 범위/길이/선택지 & 단위 \\
      \midrule
      #2
      \bottomrule
    \end{tabularx}
   } 
  \end{table}
}

모든 변수의 유효성은 필수 및 조건부 사용 여부에 따라 변수(의 속성) 값 존재의 유효성을 검사한 뒤, 타입 및 범위(또는 열거형 상수) 만족 여부를 검사하여 확인할 수 있다.

\paragraph{타입} 모든 변수는 정해진 타입의 변수여야 한다. 모든 변수의 타입은 문자\TypeTag{S}($String$), 실수\TypeTag{F}($Float$), 정수\TypeTag{I}($Integer$), 열거형\TypeTag{E}($Enum$), 불리언\TypeTag{B} 로 구분할 수 있다. 

\paragraph{필수 여부} \excel 에서는 각 변수의 필수 입력 여부에 따라, 셀 색상을 구분하고 있으며, 입력되지 않을 경우 에러가 발생하는 필수 변수 \InputTag{R}($Required$)와 달리, 비 필수 변수 \InputTag{O}($Optional$)는 미 입력 시 기본값이 사용된다. 단, 입력하지 않았다는 것을 명시적으로 \texttt{null}로 지정해야 하고, 변수 자체를 생락하면 안된다. 사용자 편의를 위한 정보 제공 목적의 변수는 사용자가 직접 입력하는 항목이 아니며, \excel 에서 \InputTag{D}로 구분된다.

\paragraph{조건부 사용 여부} 어떤 변수는 다른 변수(주로 \texttt{type})의 값에 따라 조건부로 사용된다. 이 경우에도 해당 조건에서 필수적으로 요구되는 변수 \ReqTag{CR}($Conditional Required$)와 그렇지 않은 변수 \ReqTag{CO}($Conditional Optional$)로 구분할 수 있으며, \excel 에서는 각각 \InputTag{R}과 \InputTag{D}로 표기된다.


\subsubsection{기타 규칙}

\paragraph{단위} 모든 변수는 SI 기본단위를 사용하며, \excel 에서 변수명에 명시된 단위를 적용한다 (예: $m$, $J$, $W$, $kg$).

\subsection{class별 명세}

\subsubsection{건물정보}
\label{subsubsection:ioref:building}
실재하는 건물에 대한 최상위 정보 체계이다.


\begin{table}[ht]
  \caption{건물정보 클래스의 속성들}
  \label{tbl:building_info}  
  \centering
  \begin{tabular}{lcccccc}
    \toprule
    변수명 & 타입 & 필수 & 기본값 & 조건 & 범위/길이/선택지 & 단위\\ \midrule
    건물명 & \TypeTag{S} & \InputTag{R} & - & - & - & -\\
    north\_axis & \TypeTag{F} & \InputTag{R} & - & - & [0,360) & deg(°) \\
    주소  & \TypeTag{S} & \InputTag{R} & - & - & 상세설명참고 & -\\
    허가일자 & {[ \TypeTag{I} ]} & \InputTag{R} & - & - & 상세설명참고 & -\\
    지상층수 & \TypeTag{I} & \InputTag{R} & - & - & [1,Inf) & -\\
    지하층수 & \TypeTag{I} & \InputTag{R} & - & - & [0,Inf) & -\\
    \bottomrule
  \end{tabular}
\end{table}


\paragraph{건물명} 건물의 이름. 시뮬레이션에 직접 사용되지 않는다.
\paragraph{north\_axis} 건물의 방위각. 도면상에서 나침반이 가리키는 방향 (반시계방향), 또는 건물의 향 관련 입력을 모두 회전시키는 동작 (시계방향) 을 지칭한다 (\cref{fig:ioref:north_axis}).

\begin{defaultfigure}
  
  \begin{subfigure}[b]{0.45\textwidth}
    \centering
    \includegraphics[scale=0.6]{north_axis 설명도식 (도면에서).png}
    \caption{도면에서}
  \end{subfigure}
  \hfill
  \begin{subfigure}[b]{0.45\textwidth}
    \centering
    \includegraphics[scale=0.6]{north_axis 설명도식 (지도에서).png}
    \caption{지도에서}
  \end{subfigure}
  
  \caption{north\_axis의 의미와 surface azimuth와의 관계}
  \label{fig:ioref:north_axis}

\end{defaultfigure}

\paragraph{주소} 건물의 주소. 식별 가능한 법정 기초행정구역명(세종특별자치시, 제주특별자치도 제주시 및 서귀포시 포함)을 포함하여 시/군/구 수준까지만 작성하며, 상세주소는 작성하지 않는다.
주소 정보는 \epreb 에서 허가일자와 함께 고려되어 건물의 기후지역(중부1, 중부2, 남부 등)을 판별하는 데 활용되며, 이에 따라 법적 기준 열관류율 등의 기준값을 적용하는 데 사용된다.\footnote{이에 대한 상세한 기술은 \ref{subsubsection:ioref:surface} 참고}  또한, 기상데이터(EPW)와의 연계에도 활용된다.\footnote{이에 대한 상세한 기술은 \ref{sec:address_weather_matching} 참고}

\paragraph{허가일자} 건물의 허가일자. YYYY(연, 4자리 정수)-MM(월, 2자리 정수)-DD(일, 2자리 정수)의 형식으로 입력한다. 이는, \epreb에서 법적 열관류율 등의 기준값을 적용하는 데 사용된다.\footnote{이에 대한 상세한 기술은 \ref{subsubsection:ioref:surface} 참고}

\paragraph{지상층수} 건물의 지상 층 수. 시뮬레이션에 직접 사용되지 않는다.

\paragraph{지하층수} 건물의 지하 층 수. 시뮬레이션에 직접 사용되지 않는다.

\subsubsection{실}
\label{subsubsection:ioref:zone}
건물에서 하나의 실, 또는 같은 열적 특성(유사한 용도, 공통 설비 등)을 공유하는 서로 연결된 여러 실을 모사하는 최소 단위이다.

\renewcommand{\arraystretch}{1.4}
\begin{table}[ht]
  \caption{실 클래스의 속성들}
  \label{tbl:zone_info}  
  \centering
  \begin{tabular}{lcccccc}
    \toprule
    변수명 & 타입 & 필수 & 기본값 & 조건 & 범위/길이/선택지 & 단위\\ \midrule
    층 & \TypeTag{I} & \InputTag{R} & 0 & - & - & -\\
    천정고 & \TypeTag{F} & \InputTag{R} & - & - & (0,Inf) & m \\
    이름  & \TypeTag{S} & \InputTag{R} & - & - & - & -\\
    용도프로필 & \TypeTag{E} & \InputTag{R} & - & - & 상세설명참고 & -\\
    현장조사프로필 & \TypeTag{S} & \InputTag{R} & - & - & 상세설명참고 & -\\
    조명밀도 & \TypeTag{F} & \InputTag{R} & - & - & [0,Inf) & W/m$^2$\\
    침기율 & \TypeTag{F} & \InputTag{R} & - & - & [0,Inf) & ACH50 \\
    난방 공급 설비 & \TypeTag{S} & \InputTag{R} & - & - &  {\hyperref[subsubsection:ioref:supplysystem]{\texttt{난방 공급 설비}}의 \texttt{이름}} & -\\
    냉방 공급 설비 & \TypeTag{S} & \InputTag{R} & - & - & {\hyperref[subsubsection:ioref:supplysystem]{\texttt{냉방 공급 설비}}의 \texttt{이름}} & -\\
    환기 설비 & \TypeTag{S} & \InputTag{R} & - & - & {\hyperref[subsubsection:ioref:ventilationsystem]{\texttt{환기 설비}}의 \texttt{이름}} & -\\
    면 개수 & - & \InputTag{D} & - & - & - & 개 \\
    바닥 면적 & - & \InputTag{D} & - & - & - & m$^2$\\
    \bottomrule
  \end{tabular}
\end{table}


\paragraph{층} 층을 구별하기 위한 식별자. 시뮬레이션에 직접 사용되지 않는다.

\paragraph{천정고} 존의 높이. 실의 부피 계산에 이용된다.


\paragraph{이름} 사용자 편의를 위해 제공되는 존 구별 식별자. 중복, 사용가능한 문자에 제한이 없으며, 시뮬레이션에 직접 사용되지 않는다.

\paragraph{용도프로필} 존의 용도프로필. 건축물에너지효율등급인증규정에서 제공하는 항목 중에서 선택 입력 가능하다(\cref{tbl:eco2profiles}). 

\paragraph{현장조사프로필} 현장조사프로필의 이름. \footnote{이에 대한 상세한 기술은 \ref{subsubsection:ioref:surface} 참고} @여기

\paragraph{조명밀도} 조명 밀도. 해당 실에 설치된 모든 조명기구의 정격 소비전력을 합산한 후, 이를 해당 실의 바닥면적으로 나누어 산정한다.

\paragraph{침기율} 실의 침기율. \cref{tbl:airtightness_standard}와 같은 기준으로 입력한다.

\begin{table}[h]
  \caption{준공연도별 기밀성능 기준 (n$_{50}$, h$^{-1}$)}
  \label{tbl:airtightness_standard}
  \centering
  \begin{tabular}{lccc}
    \toprule
    준공연도 & GR전 & GR후 & GR후* \\ \midrule
    1987        & 14.2 & 10.0 & 7.1 \\
    1988--2000  & 9.5  & 6.7  & 4.8 \\
    2001--2010  & 6.3  & 4.4  & 3.1 \\
    2011--2017  & 4.6  & 3.2  & 2.3 \\
    2018        & 4.0  & 2.8  & 2.0 \\
    \bottomrule
    \multicolumn{4}{l}{\footnotesize GR후: 준공연도별 약 30\% 성능 향상 (슬라이딩 창호)} \\
    \multicolumn{4}{l}{\footnotesize GR후*: 준공연도별 약 50\% 성능 향상 시스템 (이외 창호)} \\
  \end{tabular}
\end{table}


\paragraph{난방 공급 설비} 난방 공급 설비의 이름. (\ref{subsubsection:ioref:supplysystem} 참고)

\paragraph{냉방 공급 설비} 냉방 공급 설비의 이름. (\ref{subsubsection:ioref:supplysystem} 참고)

\paragraph{환기 설비} 환기시스템의 이름. (\ref{subsubsection:ioref:ventilationsystem} 참고)

\paragraph{면 개수} 해당 실에 연결된 면의 개수를 의미하며, 사용자 입력값이 아닌, 입력된 정보에 기반하여 자동으로 산정되는 값이다.

\paragraph{바닥 면적} 해당 실의 바닥 면적을 의미하며, 사용자 입력값이 아닌, 입력된 정보에 기반하여 자동으로 산정되는 값이다.


\subsubsection{면} \label{subsubsection:ioref:surface}
존과 존, 외부 (외기 또는 지면), 또는 단열선 사이의 면(벽체)이다. 하나의 실은 \texttt{실}을 인접존으로 하는 \texttt{면}을 포함하여, 1개 이상의 천장(\texttt{유형} = ``ceiling''), 1개 이상의 바닥(\texttt{type} = ``floor'') 및 1개 이상의 벽체(\texttt{type} = ``wall'')를 필요로 한다.

\renewcommand{\arraystretch}{1.4}
\begin{table}[ht]
  \caption{면 클래스의 속성들}
  \label{tbl:surface_info}  
  \centering
  \begin{tabular}{lcccccc}
    \toprule
    변수명 & 타입 & 필수 & 기본값 & 조건 & 범위/길이/선택지 & 단위\\ \midrule
    이름 & \TypeTag{S} & \InputTag{R} & - & - & - & -\\
    소속 실 & \TypeTag{S} & \InputTag{R} & - & - & - & - \\
    유형  & \TypeTag{E} & \InputTag{R} & - & - & \makecell[c]{floor \\ wall \\ ceiling} & -\\
    경계조건 & \TypeTag{E} & \InputTag{R} & - & \makecell[c]{outdoors \\ zone \\ ground \\ adiabatic} & - & -\\
    면적 & \TypeTag{F} & \InputTag{R} & - & - & (0,Inf) & m$^2$ \\
    향 & \TypeTag{F} & \ReqTag{CR} & - & \makecell[c]{\texttt{유형} = `wall' \\ AND \\ \texttt{경계조건} = `outdoors'} & [0,360) & deg(°) \\
    인접존 이름 & \TypeTag{S} & \ReqTag{CR} & - & \makecell[c]{\texttt{경계조건} = `zone'} & - & - \\
    구조체 이름 & \TypeTag{S} & \ReqTag{CR} & - & - &  \makecell[c]{\hyperref[subsubsection:ioref:construction]{\texttt{구조체}}의 이름 \\ open \\ null} & -\\
    쿨루프 반사율 & \TypeTag{F} & \ReqTag{CO} & - &
\makecell[c]{\texttt{유형} = `ceiling'\\ AND \\ \texttt{경계조건} = `outdoors'} & - & -\\
    창문 개수 & - & \InputTag{D} & - & - & - & 개 \\
    문 개수 & - & \InputTag{D} & - & - & - & 개 \\
    \bottomrule
  \end{tabular}
\end{table}


\paragraph{이름} 면 구별 식별자. 중복을 허용하지 않으며, 시뮬레이션에 직접 사용되지 않는다. 

\paragraph{소속 실} 각 면이 귀속된 실의 이름.

\paragraph{유형} 면의 유형. 천장(``ceiling''), 바닥(``floor''), 벽체(``wall'')로 구분된다. 유형이 바닥(``floor'')인 면의 면적은 해당 \texttt{실}의 면적에 포함되며, 유형이 천장(``ceiling'')인 면은 경계조건이 외기(``outdoors'')인 경우 쿨루프 적용이 가능하다.

\paragraph{경계조건} 면의 경계조건. 일사와 바람의 영향을 받고 외기와 열교환을 하는 외기(``outdoors''), 다른 \texttt{실}과 열교환을 하는 존(``zone''), 지면과 열교환을 하는 지면(``ground''), 어떠한 열교환도 하지 않고 축열의 역할만 담당하는 단열(``adiabatic'')로 구분된다 (Figure \ref{fig:surface_bo_description}).

\begin{defaultfigure}
  \includegraphics[width=0.8\textwidth]{면 경계조건 예시.png}
  \caption{면의 경계조건 예시}
  \label{fig:surface_bo_description}
\end{defaultfigure}

\paragraph{면적} 면적. 해당 면에 개구부(예:창, 문 등)가 존재하는 경우, 이들의 면적을 포함한 전체 면적을 의미한다.

\paragraph{향} 면의 방위각. 기준은 북향을 0$^\circ$로 하며, 반시계 방향으로 증가한다. 예를들어, 서향은 90$^\circ$, 남향은 180$^\circ$, 동쪽은 270$^\circ$로 정의한다.

\paragraph{인접존 이름} 인접존 이름. 면의 경계조건을 ``zone''으로 설정한 경우에만 입력한다.

\paragraph{구조체 이름} 구조체 이름. 

\paragraph{쿨루프 반사율} 지붕면의 태양 반사율. 쿨루프(Cool Roof)시공이 적용된 경우에만 입력한다.

\paragraph{창문 개수} 해당 면에 연결된 창문의 개수를 의미하며, 사용자 입력값이 아닌, 입력된 정보에 기반하여 자동으로 산정되는 값이다.

\paragraph{문 개수} 해당 면에 연결된 문의 개수를 의미하며, 사용자 입력값이 아닌, 입력된 정보에 기반하여 자동으로 산정되는 값이다.




\subsubsection{개구부} \label{subsubsection:ioref:fenestration}
개구부. 벽체의 열린 부분을 의미하며, 창, 문(유리문, 방화문 등) 등 외기와의 접촉 또는 채광·환기를 목적으로 설치된 요소를 포함한다.



\renewcommand{\arraystretch}{1.4}
\begin{table}[ht]
  \caption{개구부 클래스의 속성들}
  \label{tbl:fenestration_info}  
  \centering
  \begin{tabular}{lcccccc}
    \toprule
    변수명 & 타입 & 필수 & 기본값 & 조건 & 범위/길이/선택지 & 단위\\ \midrule
    이름 & \TypeTag{S} & \InputTag{R} & - & - & - & -\\
    소속 면 & \TypeTag{S} & \InputTag{R} & - & - & - & - \\
    유형  & \TypeTag{E} & \InputTag{R} & - & - & \makecell[c]{window \\ door \\ glassdoor} & -\\
    면적 & \TypeTag{F} & \InputTag{R} & - & - & (0,Inf) & m$^2$ \\
    구조체 이름 & \TypeTag{S} & \ReqTag{R} & - & - &  \makecell[c]{\hyperref[subsubsection:ioref:fenestrationconstruction]{\texttt{개구부 구조체}}의 \texttt{이름}} & -\\
    블라인드 & \TypeTag{E} & \ReqTag{O} & null & - & \makecell[c]{TRUE \\ FALSE} & -\\
    \bottomrule
  \end{tabular}
\end{table}

\paragraph{이름} 사용자 편의를 위해 제공되는 개구부 구별 식별자. 중복, 사용가능한 문자에 제한이 없으며, 시뮬레이션에 직접 사용되지 않는다.

\paragraph{소속 면} 각 개구부가 귀속된 면의 이름.

\paragraph{유형} 개구부의 유형.
유리창(``window''), 유리문(``glassdoor''), 또는 유리가 아닌 일반 문(``door'', 예: 방화문, 나무문 등)으로 구분된다.


\paragraph{면적} 개구부의 면적.
해당 개구부가 속한 면의 면적보다 작아야 한다.

\paragraph{구조체 이름} 개구부에 적용되는 구조체를 의미한다. 
\texttt{유형}이 ``window'' 또는 ``glassdoor''인 경우, {\hyperref[subsubsection:ioref:fenestrationconstruction]{\texttt{개구부 구조체}}}의 \texttt{투명여부} 속성값은 \texttt{투명} 이어야 한다.
\texttt{유형}이 ``door''인 경우, \texttt{투명여부} 속성값은 \texttt{불투명}이어야 한다.

\paragraph{블라인드} 블라인드 설치 여부. 값이 \texttt{null} 혹은 \texttt{FALSE}이면 블라인드가 없는 것을 의미하며, 값이 \texttt{TRUE}이면 EnergyPlus의 기본 shade가 적용된다. 



\subsubsection{공급 설비} \label{subsubsection:ioref:supplysystem}
각 존에 존재하는 말단 공급설비이다. 


\renewcommand{\arraystretch}{1.4}
\begin{table}[ht]
  \caption{공급 설비 클래스의 속성들}
  \label{tbl:supplysystem_info}  
  \centering
  \begin{tabular}{lcccccc}
    \toprule
    변수명 & 타입 & 필수 & 기본값 & 조건 & 범위/길이/선택지 & 단위\\ \midrule
    이름 & \TypeTag{S} & \InputTag{R} & - & - & - & -\\
    유형  & \TypeTag{E} & \InputTag{R} & - & - & \makecell[c]{패키지에어컨 \\ 공조기 \\ 팬코일유닛 \\ 방열기 \\ 전기방열기 \\ 바닥난방} & -\\
    냉방COP & \TypeTag{F} & \ReqTag{CR} & 3.0 & 유형 = `패키지에어컨' & (0,Inf) & W/W \\
    냉방용량 & \TypeTag{F} & \ReqTag{CR} & null & 유형 = `패키지에어컨' &  (0,Inf) & W \\
    난방용량 & \TypeTag{F} & \ReqTag{CR} & null & \makecell[c]{유형 = `방열기' \\ OR \\ `전기방열기'}  &  (0,Inf) & W \\
    생산설비명 & \TypeTag{S} & \ReqTag{CR} & - & - & \makecell[c]{\hyperref[subsubsection:ioref:sourcesystem]{\texttt{생산설비}}의 \texttt{이름}} & -\\
    \bottomrule
  \end{tabular}
\end{table}

\paragraph{이름} 사용자 편의를 위해 제공되는 공급설비 구별 식별자. 중복, 사용가능한 문자에 제한이 없으며, 시뮬레이션에 직접 사용되지 않는다.

\paragraph{유형} 공급설비의 종류. 선택한 공급설비 유형에 따라 냉·난방 제공 가능 여부가 달라진다. 각 공급설비별 냉·난방 지원 가능 여부는 Figure \ref{fig:supply_and_source_properties_and_connection}에 제시되어 있다.

\paragraph{냉방COP} 정격 냉방 성능계수(COP). 정격 운전 조건에서의 냉방 출력을 소비 전력으로 나눈 값이다. 공급설비의 유형을 ``패키지에어컨''으로 선택한 경우에만 입력한다.

\paragraph{냉방용량} 설비의 최대 냉방 출력을 의미한다. 해당 값을 입력하지 않을 경우 (\texttt{null}), EnergyPlus의 autosizing을 이용한다.\footnote{Autosizing은 설계일(design day) 조건에서의 피크 냉방부하를 기준으로, 
설비가 해당 부하를 만족할 수 있도록 용량을 자동으로 산정하는 방법이다. 
냉방 설계일은 일반적으로 ASHRAE 기준에 따른 연간 0.4\% 빈도 조건의 고온 기상 데이터를 사용한다.} 공급설비의 유형을 ``패키지에어컨''로 선택한 경우에만 입력한다.


\paragraph{난방용량} 설비의 최대 난방 출력을 의미한다. 해당 값을 입력하지 않을 경우 (\texttt{null}), EnergyPlus의 autosizing을 이용한다. 공급설비의 유형을 ``방열기''혹은 ``전기 방열기''로 선택한 경우에만 입력한다.

\paragraph{생산설비} 공급설비와 연결된 생산설비의 이름. 연결가능한 생산설비는 \cref{fig:supply_and_source_properties_and_connection}\와 같다.

\begin{defaultfigure}
  \includegraphics[width=\textwidth]{HVAC_properties_and_connectables.png}
  \caption{공급 및 생산설비 종류에 따른 요구 속성 및 연결 가능성 (적색 선은 난방, 청색 선은 냉방이 가능함을 의미하며, packaged\_air\_conditioner와 electric\_radiator는 생산설비의 연결을 요구하지 않는다.)}
  \label{fig:supply_and_source_properties_and_connection}
\end{defaultfigure}


\subsubsection{생산설비} \label{subsubsection:ioref:sourcesystem}
건물에 존재하는 생산설비이다.


\begin{longtable}{lcccccc}
\caption{생산 설비 클래스의 속성들} \label{tbl:sourcesystem_info} \\

\toprule
변수명 & 타입 & 필수 & 기본값 & 조건 & 범위/길이/선택지 & 단위\\ 
\midrule
\endfirsthead

\toprule
변수명 & 타입 & 필수 & 기본값 & 조건 & 범위/길이/선택지 & 단위\\ 
\midrule
\endhead

\midrule
\multicolumn{7}{r}{다음 페이지 계속} \\
\endfoot

\bottomrule
\endlastfoot

이름 & \TypeTag{S} & \InputTag{R} & - & - & - & -\\
유형  & \TypeTag{E} & \InputTag{R} & - & - & \makecell[c]{히트펌프 \\ 지열히트펌프 \\ 냉동기 \\ 흡수식냉동기 \\ 지역난방 \\ 보일러 } & -\\
냉방COP & \TypeTag{F} & \ReqTag{CO} & 3.0 & \makecell[c]{유형 = `냉동기' \\ OR \\ `히트펌프' \\ OR \\ `지열히트펌프' \\ OR \\ `흡수식냉동기'} & (0,Inf) & W/W \\
냉방용량 & \TypeTag{F} & \ReqTag{CO} & null & \makecell[c]{유형 = `냉동기' \\ OR \\ `히트펌프' \\ OR \\ `지열히트펌프' \\ OR \\ `흡수식냉동기'} & (0,Inf) & W \\
난방COP & \TypeTag{F} & \ReqTag{CO} & 3.0 & \makecell[c]{유형 = `히트펌프' \\ OR \\ `지열히트펌프'} & (0,Inf) & W/W \\
난방용량 & \TypeTag{F} & \ReqTag{CR} & null & \makecell[c]{유형 = `히트펌프' \\ OR \\ `지열히트펌프' \\ OR \\ `보일러'} & (0,Inf) & W \\
효율 & \TypeTag{F} & \ReqTag{CO} & 0.85 & 유형 = `보일러' & (0,1) & - \\
급탕용 & \TypeTag{B} & \ReqTag{CR} & - & \makecell[c]{유형 = `보일러' \\ OR \\ 지역난방} & (0,Inf) & - \\
연료종류 & \TypeTag{E} & \ReqTag{CR} & 3.0 & 유형 $\neq$ `지역난방' & \makecell[c]{전기 \\ 가스 \\ 난방유 \\ 지역난방} & - \\
압축기 종류 & \TypeTag{E} & \InputTag{R} & - & 유형 = `냉동기' & \makecell[c]{터보식 \\ 스크류식 \\ 왕복동식} & - \\
냉각탑 종류 & \TypeTag{E} & \ReqTag{CR} & - & 유형 = `냉동기' & \makecell[c]{개방형 \\ 밀폐형} & - \\
냉각탑 용량 & \TypeTag{F} & \ReqTag{CO} & null & 유형 = `냉동기' & (0,Inf) & W \\
냉각탑 제어방식 & \TypeTag{E} & \ReqTag{CR} & - & 유형 = `냉동기' & \makecell[c]{single-speed \\ two-speed} & - \\
연동보일러 효율 & \TypeTag{F} & \ReqTag{CO} & 0.85 & 유형 = `흡수식냉동기' & (0,1) & - \\

\end{longtable}



\paragraph{이름} 사용자 편의를 위해 제공되는 생산설비 구별 식별자. 중복, 사용가능한 문자에 제한이 없으며, 시뮬레이션에 직접 사용되지 않는다.

\paragraph{유형} 생산설비의 종류. 공급설비에 열을 제공하는 기기의 종류를 선택하며, 선택가능한 유형은 히트펌프(``heatpump''), 지열 히트펌프(``geothermal\_heatpump''), 칠러(``chiller''), 흡수식 칠러(``absorption\_chiller''), 보일러(``boiler''), 지역난방(``district\_heating'')이다.

\paragraph{냉방COP} 정격 냉방성능계수(COP). 정격 운전 조건에서의 냉방 출력을 소비 전력으로 나눈 값이다.

\paragraph{냉방용량} 해당 값을 입력하지 않을 경우 (\texttt{null}) EnergyPlus의 autosizing을 이용한다.

\paragraph{난방COP} 정격 난방성능계수(COP). 정격 운전 조건에서의 난방 출력을 소비 전력으로 나눈 값이다.

\paragraph{난방용량} 해당 값을 입력하지 않을 경우 (\texttt{null}) EnergyPlus의 autosizing을 이용한다.

\paragraph{효율} 보일러의 효율.

\paragraph{급탕용} 급탕공급 여부. 급탕을 공급하는 모든 생산설비들은 급탕효율 계산에 이용되며, 모든 설비가 동등하게 급탕 부하를 분담하는 것으로 가정한다.

\paragraph{연료종류} 연료의 종류. 생산설비의 사용 연료를 의미하며, 선택 가능한 항목은 전기(``electricity''), 천연가스/가스(``natural\_gas''), 난방유(``oil''), 지역난방(``district\_heating'')이다.

\paragraph{압축기 종류} 냉동기의 냉매를 압축하는 방식에 따른 압축기의 유형을 의미한다. 
원심력을 이용하는 \texttt{터보식}, 나사형 로터의 회전에 의해 압축하는 \texttt{스크류식}, 
피스톤의 왕복 운동에 의해 압축하는 \texttt{왕복동식} 중 선택하여 입력한다.

\paragraph{냉각탑 종류} 냉각탑의 열교환 방식에 따른 유형을 의미한다. 
냉각수가 외기와 직접 접촉하는 \texttt{개방형}과, 
열교환기를 통해 간접적으로 열을 교환하는 \texttt{밀폐형} 중 선택하여 입력한다.

\paragraph{냉각탑 용량} 해당 값을 입력하지 않을 경우 (\texttt{null}) EnergyPlus의 autosizing을 이용한다.

\paragraph{냉각탑 제어방식} 냉각탑에 적용된 제어. 

\paragraph{연동보일러 효율} 연동된 보일러의 효율.

\subsubsection{환기설비} \label{subsubsection:ioref:ventilationsystem}


\renewcommand{\arraystretch}{1.2}
\begin{table}[h]
  \caption{환기설비 클래스의 속성들}
  \label{tbl:ventilationsystem_info}  
  \centering
  \begin{tabular}{lcccccc}
    \toprule
    변수명 & 타입 & 필수 & 기본값 & 조건 & 범위/길이/선택지 & 단위\\ \midrule
    이름 & \TypeTag{S} & \InputTag{R} & - & - & - & -\\
    난방효율  & \TypeTag{F} & \InputTag{O} & 0.70 & - & (0,1) & - \\
    냉방효율 & \TypeTag{F} & \InputTag{O} & 0.45 & - & (0,1) & - \\
    \bottomrule
  \end{tabular}
\end{table}

\paragraph{이름} 사용자 편의를 위해 제공되는 생산설비 구별 식별자. 중복, 사용가능한 문자에 제한이 없으며, 시뮬레이션에 직접 사용되지 않는다.

\paragraph{난방효율} 유효 냉방 전열교환 효율.

\paragraph{냉방효율} 유효 난방 전열교환 효율.


\subsubsection{PV패널} \label{subsubsection:ioref:photovoltaicsystem}
태양광 발전 설비.

\renewcommand{\arraystretch}{1.2}
\begin{table}[h]
  \caption{PV패널 클래스의 속성들}
  \label{tbl:photovoltaicsystem_info}  
  \centering
  \begin{tabular}{lcccccc}
    \toprule
    변수명 & 타입 & 필수 & 기본값 & 조건 & 범위/길이/선택지 & 단위\\ \midrule
    이름 & \TypeTag{S} & \InputTag{R} & - & - & - & -\\
    면적  & \TypeTag{F} & \InputTag{R} & - & - & (0,Inf) & m$^2$ \\
    효율 & \TypeTag{F} & \InputTag{R} & - & - & (0,1) & - \\
    방위각 & \TypeTag{F} & \InputTag{R} & - & - & [0,360) & deg(°) \\
    경사각 & \TypeTag{F} & \InputTag{R} & - & - & [0,90] & deg(°) \\
    \bottomrule
  \end{tabular}
\end{table}


\paragraph{이름} 사용자 편의를 위해 제공되는 태양광설비 구별 식별자. 중복, 사용가능한 문자에 제한이 없으며, 시뮬레이션에 직접 사용되지 않는다.

\paragraph{면적} 태양광 패널의 면적.

\paragraph{효율} 태양광 모듈의 발전 효율.

\paragraph{방위각} 패널의 방위각. \hyperref[subsubsection:ioref:building]{\texttt{building}}의 \texttt{north\_axis}를 고려하지 않고 독립적으로 계산된다 (Figure \ref{fig:azimuth_of_PVpanel}). 

\begin{defaultfigure}
  \includegraphics{태양광패널의 azimuth 설명.png}
  \caption{태양광 패널의 \texttt{azimuth}}
  \label{fig:azimuth_of_PVpanel}
\end{defaultfigure}

\paragraph{경사각} 패널의 경사각. 0$^\circ$는 수평면, 90$^\circ$는 수직면을 의미한다.


\subsubsection{재료} \label{subsubsection:ioref:material}

\renewcommand{\arraystretch}{1.2}
\begin{table}[h]
  \caption{재료 클래스의 속성들}
  \label{tbl:material_info}  
  \centering
  \begin{tabular}{lcccccc}
    \toprule
    변수명 & 타입 & 필수 & 기본값 & 조건 & 범위/길이/선택지 & 단위\\ \midrule
    이름 & \TypeTag{S} & \InputTag{R} & - & - & - & -\\
    열전도율  & \TypeTag{F} & \InputTag{R} & - & - & - & W/m·K \\
    밀도 & \TypeTag{F} & \InputTag{R} & - & - & - & kg/m$^3$ \\
    비열 & \TypeTag{F} & \InputTag{R} & - & - & (100,Inf) & J/kg·K \\
    \bottomrule
  \end{tabular}
\end{table}

\paragraph{이름} 사용자 편의를 위해 제공되는 재료 구별 식별자. 중복, 사용가능한 문자에 제한이 없으며, 시뮬레이션에 직접 사용되지 않는다.

\paragraph{열전도율} 재료의 단위 두께를 통한 열 전달 능력을 나타낸다.

\paragraph{밀도} 재료의 단위 부피당 질량을 나타낸다.

\paragraph{비열} 재료의 단위 질량이 1K 상승할 때 필요한 열량을 나타낸다.

\subsubsection{면 구조체} \label{subsubsection:ioref:construction}

\renewcommand{\arraystretch}{1.2}
\begin{table}[h]
  \caption{면 구조체 클래스의 속성들}
  \label{tbl:surface_construction_info}  
  \centering
  \begin{tabular}{lcccccc}
    \toprule
    변수명 & 타입 & 필수 & 기본값 & 조건 & 범위/길이/선택지 & 단위\\ \midrule
    이름 & \TypeTag{S} & \InputTag{R} & - & - & - & -\\
    레이어 재료  & \TypeTag{S} & \InputTag{R} & - & - & \hyperref[subsubsection:ioref:material]{\texttt{재료}}의 \texttt{이름} & - \\
    레이어 두께  & \TypeTag{F} & \InputTag{R} & - & - & (0.003, Inf) & mm \\
    \bottomrule
  \end{tabular}
\end{table}

\paragraph{이름} 사용자 편의를 위해 제공되는 레이어 구별 식별자. 중복, 사용가능한 문자에 제한이 없으며, 시뮬레이션에 직접 사용되지 않는다.

\paragraph{레이어 재료} 레이어 구성 재료의 이름.

\paragraph{레이어 두께} 레이어 구성 재료의 두께.


\subsubsection{개구부 구조체} \label{subsubsection:ioref:fenestrationconstruction}

\renewcommand{\arraystretch}{1.2}
\begin{table}[h]
  \caption{개구부 구조체 클래스의 속성들}
  \label{tbl:fenestration_construction_info}  
  \centering
  \begin{tabular}{lcccccc}
    \toprule
    변수명 & 타입 & 필수 & 기본값 & 조건 & 범위/길이/선택지 & 단위\\ \midrule
    이름 & \TypeTag{S} & \InputTag{R} & - & - & - & -\\
    투명여부  & \TypeTag{B} & \InputTag{R} & - & - & \makecell[c]{TRUE \\ FALSE} & - \\
    열관류율  & \TypeTag{F} & \InputTag{R} & - & - & - & W/m$^2$·K \\
    태양열취득계수  & \TypeTag{F} & \ReqTag{CR} & - & 투명여부 = `TRUE' & (0,1] & - \\
    \bottomrule
  \end{tabular}
\end{table}

\paragraph{이름} 창호 객체의 이름. 중복을 허용하지 않으며, 시뮬레이션에 직접 사용되지 않는다.

\paragraph{투명여부} 구조체의 투명여부. 투명한 구조체의 개구부(예: 유리창, 유리문 등)는 ``TRUE''로 입력하며, 불투명한 구조체의 개구부(예: 방화문, 나무문 등)는 ``FALSE''로 입력한다.

\paragraph{u} 창호의 열관류율(U-value).

\paragraph{g} 창호의 열취득계수(G-value, SHGC).

% ---------------------------------------------------------------------------- %
%                                  NEW SECTION                                 %
% ---------------------------------------------------------------------------- %

\section{입력파일 (xlsx)}
\label{section:format-explanation-xlsx}

엑셀 기반 입력파일은 각 시트가 grm 입력파일의 클래스에 대응되는 형태로 구성된다. 각 시트 내에서, 열은 클래스의 속성에 대응되며, 행은 개별 객체 인스턴스에 해당한다.


% ---------------------------------------------------------------------------- %
%                                  NEW SECTION                                 %
% ---------------------------------------------------------------------------- %

\section{출력파일 (grr)}
\label{section:format-explanation-grr}

GRR 파일은 시뮬레이션 수행 결과를 저장하기 위한 데이터 구조로, 건물의 에너지 사용량 및 이를 기반으로 계산된 다양한 성능 지표를 포함한다. GRM 파일이 건물 모델과 시뮬레이션 입력 정보를 정의하는 역할을 하는 반면, GRR 파일은 시뮬레이션 결과 데이터를 체계적으로 저장하고 후속 분석이나 보고서 생성에 활용하기 위한 목적으로 생성된다. 데이터 형식은 JSON 기반 구조를 따르며, 건물 기본 정보, 계산에 사용되는 변환계수, 용도별 에너지 사용량, 그리고 집계된 성능 지표로 구성된다.

GRR 데이터의 최상위 구조는 다음과 같은 주요 객체로 구성된다.

\begin{itemize}
\item \texttt{building} : 건물의 기본 정보
\item \texttt{constants} : 에너지 변환계수 및 비용계수
\item \texttt{site\_uses} : 용도별 연료 에너지 사용량
\item \texttt{source\_uses} : 1차 에너지 기준 사용량
\item \texttt{co2} : 온실가스 배출량
\item \texttt{cost} : 에너지 비용
\item \texttt{summary\_per\_area} : 면적당 성능 지표
\item \texttt{summary\_gross} : 건물 전체 기준 성능 지표
\end{itemize}

\subsection{class별 명세}

\subsubsection{건물 정보}

\texttt{building} 객체는 결과 계산에 사용된 바닥면적을 포함하여 건물의 기본 정보를 저장한다.

\subsubsection{변환계수 (constants)}

에너지 사용량을 1차 에너지, 온실가스 배출량, 그리고 비용으로 변환하기 위해 사용되는 계수들이 \texttt{constants} 객체에 정의된다.

\begin{verbatim}
constants
 ├─ site2source
 ├─ site2co2
 └─ site2cost
\end{verbatim}

각 항목은 연료 유형(전력:\texttt{electricity}, 천연가스:\texttt{natural\_gas}, 지역난방:\texttt{district\_heating}, 유류:\texttt{oil})별 변환계수를 정의한다.

\begin{itemize}
\item \texttt{site2source} : 에너지소요량[$kWh$]를 1차에너지소요량[$kWh$]로 변환하기 위한 계수
\item \texttt{site2co2} : 에너지소요량[$kWh$]을 온실가스 배출량[$kgCO\textsubscript{2,eq}$]으로 변환하기 위한 계수
\item \texttt{site2cost} : 에너지소요량[$kWh$]을 비용[$KRW$]으로 변환하기 위한 계수
\end{itemize}

\subsubsection{에너지소요량 (site\_uses [$kWh$])}

에너지소요량은 먼저 용도(end-use) 기준으로 분류되며, 각 용도는 다시 연료 유형별 월별 사용량 배열로 저장된다.

\begin{verbatim}
site_uses
 ├─ heating
     ├─ ELECTRICITY     : [12개월 배열]
     ├─ NATURALGAS      : [12개월 배열]
     ├─ OIL             : [12개월 배열]
     └─ DISTRICTHEATING : [12개월 배열]
       (다른 end-use도 동일)
 ├─ cooling
 ├─ lighting
 ├─ circulation
 ├─ hotwater
 └─ generators
\end{verbatim}

\subsubsection{1차에너지소요량 (source\_uses [$kWh$])}
구조는 에너지소요량과 같다.

\subsubsection{온실가스 배출량 (co2 [$kgCO\textsubscript{2,eq}$])}
구조는 에너지소요량과 같다.

\subsubsection{에너지 비용 (cost [$KRW$])}
구조는 에너지소요량과 같다.

\subsubsection{집계: 면적당 (summary\_per\_area)}

각 항목별 합계를 건물의 바닥면적으로 나눈 값이다.

\begin{verbatim}
summary_per_area
 ├─ site_uses
     (연료별 집계)
     ├─ ELECTRICITY     
     ├─ NATURALGAS      
     ├─ OIL             
     └─ DISTRICTHEATING 
     (end-use별 집계)
     ├─ heating     
     ├─ cooling      
     ├─ lighting
     ├─ circulation
     ├─ hotwater
     └─ generators      
     (합계)
     ├─ total_monthly : [12개월 배열]
     └─ total_annual
 ├─ source_uses
 ├─ co2
 └─ cost
\end{verbatim}

\subsubsection{집계: 총면적 (summary\_gross)}

구조는 면적당 집계와 같다.
